% 源工作目录：../../raw/videos/朱邦华- SGLang，强化学习，英伟达收购，二次创业，清华，伯克利，LMSYS，Chatbot Arena，勇于放弃/，如需重新编译请在该目录下执行 xelatex

\documentclass[a4paper]{article}

% Basic packages
\usepackage[fontset=fandol]{ctex}
\usepackage{amsmath, amssymb}
\usepackage{graphicx}
\usepackage[margin=2.5cm]{geometry}
\usepackage[most]{tcolorbox}
\usepackage{etoolbox}
\usepackage{listings}
\usepackage{hyperref}
\usepackage{booktabs}
\usepackage{subcaption}
\usepackage{float}
\usepackage{tikz}
\usepackage{xcolor}

% Highlight boxes
\newtcolorbox{knowledgebox}[1]{
    enhanced, colback=blue!5!white, colframe=blue!75!black, colbacktitle=blue!75!black,
    coltitle=white, fonttitle=\bfseries, title=#1, attach boxed title to top left={yshift=-2mm, xshift=2mm},
    boxrule=1pt, sharp corners
}

\newtcolorbox{importantbox}[1]{
    enhanced, colback=yellow!10!white, colframe=yellow!80!black, colbacktitle=yellow!80!black,
    coltitle=black, fonttitle=\bfseries, title=#1, sharp corners
}

\newtcolorbox{warningbox}[1]{
    enhanced, colback=red!5!white, colframe=red!75!black, colbacktitle=red!75!black,
    coltitle=white, fonttitle=\bfseries, title=#1, sharp corners
}

% Listings
\lstset{
    language=Python,
    basicstyle=\ttfamily\small,
    keywordstyle=\color{blue},
    stringstyle=\color{red!60!black},
    commentstyle=\color{green!60!black},
    breaklines=true,
    frame=single,
    numbers=left,
    numberstyle=\tiny\color{gray},
    captionpos=b,
    extendedchars=false
}

% Front-page metadata
\newcommand{\notetitle}{访谈笔记：朱邦华 —— SGLang、强化学习、英伟达收购与二次创业}
\newcommand{\noteauthors}{月球大叔 \& Codex}
\newcommand{\notedate}{\today}
\newcommand{\videochannel}{月球大叔}
\newcommand{\videopublishdate}{2026年7月}
\newcommand{\videoduration}{1 小时 47 分钟}
\newcommand{\videourl}{https://www.bilibili.com/video/BV1TTLc6HEKT/}
\newcommand{\videocoverpath}{cover.jpg}

\begin{document}

\begin{titlepage}
\centering
{\Large 访谈笔记\par}
\vspace{1.2cm}
{\huge\bfseries \notetitle\par}
\vspace{0.8cm}
{\large \noteauthors\par}
\vspace{0.3cm}
{\large \notedate\par}
\vspace{1.2cm}

\includegraphics[width=0.82\textwidth,height=0.45\textheight,keepaspectratio]{\videocoverpath}\par

\vfill
\begin{tcolorbox}[width=0.9\textwidth, colback=black!2!white, colframe=black!60, sharp corners]
\textbf{视频作者/频道}：\videochannel\par
\textbf{发布日期}：\videopublishdate\par
\textbf{视频时长}：\videoduration\par
\textbf{视频链接}：
\href{\videourl}{\nolinkurl{\videourl}}
\end{tcolorbox}
\end{titlepage}

\tableofcontents
\newpage

%% --- 正文内容开始 --- %%

\section{引言}

2026年7月，播客节目「月球大叔」邀请到 RadixArk 联合创始人兼 CTO 朱邦华进行了一场近两小时的深度访谈。朱邦华的经历在 AI 圈内几乎是教科书级别的传奇：清华电子系年级第一，伯克利师从 Michael Jordan 和 Jiantao Jiao 从事机器学习理论研究，博士期间预判大模型浪潮并果断从理论转向系统，创办 NexusFlow 成为开源社区最早跑通 PPO 训练的团队，随即被英伟达收购。然而在被收购不到半年后，他又毅然放弃后续的归属权益，与 SGLang 联合创始人盛颖共同创立 RadixArk，一举拿下 1 亿美元融资。

\begin{figure}[H]
\centering
\includegraphics[width=0.75\textwidth]{figures/talk_opening.jpg}
\caption{访谈开场：朱邦华与月球大叔对谈 \protect\footnotemark}
\end{figure}
\footnotetext{视频画面时间区间：00:01:30--00:02:15。}

这场对话横跨技术、创业、学术、哲学等多个维度，信息密度极高。本笔记以主题为线索重新组织访谈内容，保留朱邦华的核心观点和独特视角，同时加入必要的技术背景说明。

\subsection{本章小结}
朱邦华的人生轨迹体现了一种罕见的判断力：从理论到系统、从学术到创业、从被收购到再次出发，每一步都踩在了 AI 浪潮的关键节点上。本笔记将围绕 RadixArk 的技术使命、强化学习的实战经验、创业决策哲学、以及他对 AI 未来的判断展开。

\section{RadixArk：AI 生命周期的全栈平台}

\subsection{公司的核心定位}

朱邦华将 RadixArk 定位为一家 AI 基础设施建设公司，核心愿景是覆盖 AI 的全生命周期（full lifecycle）：

\begin{itemize}
\item \textbf{Pretraining（预训练）}：投入成本巨大，能做预训练的公司不多，更多是头部 Lab
\item \textbf{Post-training（后训练）}：需求最大、增长最快的环节，尤其是 RLHF/RLVR 方向；成本相比预训练更低，但对私有数据和垂直领域数据的需求很高
\item \textbf{Inference（推理）}：SGLang 作为推理引擎的基础，与 DeepSeek、Thinking Machines 等有深度合作
\end{itemize}

\begin{importantbox}{AI 生命周期的商业闭环}
朱邦华的判断是：从降本增效的角度，企业更愿意先通过 post-training 把模型训练到满足自身需求的状态，再去做推理部署。RadixArk 的商业逻辑是：帮助客户提升私有模型的能力，客户再使用 SGLang 做推理——形成「训练提升模型能力 $\rightarrow$ 用 SGLang 做推理部署」的完整闭环。
\end{importantbox}

\begin{knowledgebox}{SGLang 的历史渊源}
SGLang 由盛颖（Ying Sheng）和郑怜悯（Lianmin Zheng）在伯克利期间共同创建，最初是 LMSYS Org 下的一个开源推理框架项目。朱邦华在 PhD 第 4--5 年期间读到他们早期的论文《FlexGen》（一篇关于在 CPU 上做 LLM 推理的早期工作），主动联系了盛颖，三人由此开始了长达数年的研究合作。社区逐渐壮大后，他们决定不再局限于推理，而是向训练等更前沿的 Infra 方向拓展——这也正是朱邦华加入的契机。
\end{knowledgebox}

\subsection{从降本增效到能力提升}

朱邦华特别强调了 post-training 的战略意义。他以 Cursor 为例：Cursor 会先对模型（如 KIMI）进行 post-training，提升其在 Cursor 内部作为 Agent 的表现，然后再做推理部署。「帮客户提升私有模型的能力，客户再用我们的 Infra 做部署」——这就是 RadixArk 的核心商业逻辑。

\subsection{本章小结}
RadixArk 不满足于做一个纯粹的推理引擎公司，而是要成为 AI 全生命周期的 Infra 平台，尤其聚焦在 post-training 这个当前最大的痛点上。公司的差异化在于：开源社区出身的技术深度 + 与头部模型公司的深度合作 + 覆盖训练到推理的完整产品矩阵。

\section{SGLang 推理引擎的技术前沿}

\subsection{新模型架构对推理 Infra 的挑战}

朱邦华指出，当 DeepSeek V4 这样的新模型出现时，其模型架构创新会给推理 Infra 带来全新的挑战。以 DeepSeek V4 为例：

\begin{itemize}
\item \textbf{Shadow Redux}：针对 DeepSeek V4 的 prefix caching 优化方案
\item \textbf{FlashAttention 变体}：需要针对新的稀疏注意力模式做适配
\item \textbf{架构感知的调度}：不能假设所有模型共用同一套推理优化策略
\end{itemize}

\begin{importantbox}{推理引擎的核心竞争力：架构适配}
朱邦华强调，推理 Infra 不是一套「通用方案打天下」的生意。每当前沿模型在架构上做出创新，推理引擎就需要跟随适配——那些能最快适配新架构的引擎，就能吃到第一波部署红利。这也是为什么 SGLang 团队投入大量精力与 DeepSeek 等前沿 Lab 紧密合作。
\end{importantbox}

\subsection{从 FlashInfer 到 SGLang 的演进}

FlashInfer 最初是 SGLang 生态中的高性能注意力算子库，专注于为各种注意力模式提供高效 CUDA kernel。随着模型架构越来越多样化（MLA、GQA、滑动窗口注意力、稀疏注意力等），FlashInfer 的作用从「算子优化」升级为「架构适配层」——它使得 SGLang 能够在不重写推理引擎的前提下，快速支持新出现的注意力模式。

\begin{figure}[H]
\centering
\begin{tikzpicture}[node distance=0.8cm, auto]
  \tikzstyle{layer}=[rectangle, draw, minimum width=10cm, minimum height=0.7cm,
                     align=center, font=\small, rounded corners=1pt]
  \node [layer, fill=blue!10] (api)    {API / 服务层 (OpenAI-compatible)};
  \node [layer, fill=green!10, below of=api] (sched) {调度器 (RadixAttention Prefix Cache + Continuous Batching)};
  \node [layer, fill=yellow!10, below of=sched] (vm) {内存管理 (PagedAttention / KVCache Pool)};
  \node [layer, fill=orange!10, below of=vm] (kernels) {算子层 (FlashInfer / FlashAttention / Shadow Redux)};
  \node [layer, fill=red!10, below of=kernels] (backend) {后端适配 (CUDA / ROCm / x86)};
\end{tikzpicture}
\caption{SGLang 推理引擎的架构层级}
\end{figure}

\subsection{本章小结}
推理 Infra 的核心挑战已从「通用性能优化」转向「架构感知的快速适配」。SGLang 通过模块化的架构设计（调度器 + 内存管理 + FlashInfer 算子层 + 后端适配层）实现了灵活性，能够在 DeepSeek V4 等新架构出现时快速响应。这种「适配速度」本身就是一种竞争壁垒。

\section{从理论到系统：学术与创业之路}

\subsection{清华岁月：考试型天才的底色}

朱邦华来自江苏连云港（苏北），自称「苏北比较落后，一个学校能考上清华的不是特别多」。在清华电子系，他几乎每学期都是系里第一。当被问及成功的秘诀时，他给出了一个极具「大模型风格」的自嘲：

\begin{quote}
「如果用大语言模型来类比，我的 long context 能力比较强……我属于那种第二天要考固体物理，前面基本没怎么学，然后通宵看一遍书，第二天去考也能考第一的人。」\par
\hfill —— 朱邦华，00:52:00--00:53:20
\end{quote}

他不认为自己属于特别努力或智力超常的类型（「参加各种竞赛也都不太行」），但具备一种快速吸收并整合大量信息的能力——这恰是后来他在 AI 研究中展现的核心特质。

\subsection{伯克利：Michael Jordan 的熏陶}

在伯克利读博期间，朱邦华师从 Jiantao Jiao 和 Michael Jordan。他的研究领域是机器学习理论和统计学，具体方向包括：

\begin{itemize}
\item \textbf{RL Theory}：强化学习的理论分析
\item \textbf{Robust Statistics}：鲁棒统计学习，与 Jacob Steinhardt 合作写了一篇近 200 页的论文，投到 Annals of Statistics
\item \textbf{Optimization \& Economics}：在 Michael Jordan 的引导下，大量阅读优化理论、博弈论、拍卖理论等——因为 Jordan 认为「Human-AI interaction 本质上是一个经济学问题」
\end{itemize}

\begin{knowledgebox}{Michael Jordan 的跨学科视野}
朱邦华回忆，Michael Jordan 在当时就预见到 Human-AI Interaction 的重要性。他认为机器与机器的交互和机器与人的交互有本质不同——后者不再是纯技术问题，而是人如何与机器协作的经济学问题。因此他带领学生读了好几年的经济学经典，包括博弈论和拍卖理论。这种跨学科训练深刻影响了朱邦华后来的判断力。
\end{knowledgebox}

\subsection{GPT 时刻：从理论到系统的转向}

博士第五年，朱邦华第一次使用 ChatGPT。那一刻改变了他的人生轨迹：

\begin{quote}
「在这之前，我一直在做 ML theory 和 statistics、RL theory。但用 ChatGPT 的一刹那，我觉得这可能会是一个划时代的东西。」\par
\hfill —— 朱邦华，00:08:30--00:09:15
\end{quote}

他当时的判断逻辑非常清晰：在没有经过精细训练的情况下，ChatGPT 已经展现出了惊人的能力——如果继续 scale（无论是参数量、数据量还是 test-time thinking），潜力将是巨大的。Michael Jordan 最初持怀疑态度（「在统计上就是个 next-token prediction problem，和以前的东西没有任何区别」），但朱邦华坚信自己的直觉。

\begin{importantbox}{「判断力」的第一次大考}
朱邦华在 2023 年初的判断可以概括为：GPT 的表现说明 scaling 这个方向是对的，越 scale 能力越强。不是炒作，是真正的范式变化。这个判断让他果断从理论转向系统和 LLM，也最终导致了他后来创办 NexusFlow、被英伟达收购、以及今天的 RadixArk。
\end{importantbox}

\subsection{本章小结}
朱邦华的学术经历为他后来的判断力奠定了三个基础：（1）清华的「long context」能力——快速吸收大量信息并整合；（2）伯克利的理论训练——从 RL theory 到经济学，跨学科思维；（3）GPT 时刻的判断——在所有人还在争论时，他选择了 all in。这三个要素叠加，造就了他后来的职业轨迹。


\section{强化学习实战：从 PPO 到 RLVR}

\subsection{开源社区第一个跑通 PPO 的故事}

2023 年，当大多数开源团队还在做 SFT（Supervised Fine-Tuning）时，朱邦华带领的 NexusFlow 团队成为了开源社区中第一个成功用 PPO 训练 LLM 的团队。这个成就的含金量在于：PPO 训练涉及 Actor、Critic、Reward Model 三个模型的协同优化，工程复杂度远超 SFT。

\begin{quote}
「我觉得我们可能是当时唯一一家把 PPO 在开源这边调通的地方。」\par
\hfill —— 朱邦华，00:28:45--00:29:05
\end{quote}

他们当时使用开源偏好数据，做了一个 RLHF 的开源版本，模型在 Chatbot Arena 上的分数非常高。这个工作证明了：通过 RLHF，可以让模型在对话质量上实现质的飞跃。

\subsection{Critic 热启动：调通 PPO 的核心 Trick}

当被追问「调通 PPO 最重要的因素是什么」时，朱邦华给出了一个非常具体的例子——Critic 的初始化策略：

\begin{importantbox}{Critic Warmup Trick}
如果直接随机初始化 Critic（价值函数网络），训练会极不稳定。因为 PPO 的搜索空间维度极高，随机初始化的 Critic 无法提供有效的 advantage 估计。\par
\textbf{正确做法}：从 Reward Model 的权重初始化 Critic，让 Critic 一开始就具备评估响应质量的「直觉」。这个简单的初始化策略直接决定 PPO 能否收敛。
\end{importantbox}

朱邦华用了一个很好的比喻：「PPO 不是 fancy 的东西——infra 不是说起来那么简单，因为这个问题 dimension 太高，你没有办法做 grid search。你需要用自己的直觉去探索出一条路来。」

调通一个 7B 模型的 PPO 训练大约花了一个月时间，但准备数据的时间更久——因为当时开源社区几乎没有多轮偏好对比数据（一个问题对应多个不同质量的回答，并进行 ranking），NexusFlow 花了很多精力构建和开源了这样一个数据集。

\subsection{PPO vs GRPO：简单与上限的权衡}

朱邦华对 PPO 和 GRPO 的评价非常客观：

\begin{itemize}
\item \textbf{PPO}：上限更高，但工程调参难度大。「如果你把 PPO 做得好，上限是很高的。」
\item \textbf{GRPO}：相对更简单，调参更容易，现在使用的人越来越多。但本质上放弃了 PPO 通过 Critic 带来的精细优势估计。
\end{itemize}

\begin{warningbox}{PPO 的工程陷阱}
PPO 的训练链路长（Actor $\rightarrow$ Reward Model $\rightarrow$ Critic $\rightarrow$ Advantage $\rightarrow$ Policy Update），每个环节都可能出错。朱邦华强调：「它有很多地方是 implementation 和真正串起来的时候出错的地方。」没有经验的团队很容易在某个环节卡住，然后错误地认为「PPO 不 work」，实际上只是实现有 bug。
\end{warningbox}

\subsection{RLHF vs RLVR：偏好对齐与可验证奖励}

朱邦华对强化学习的分类非常清晰：

\begin{itemize}
\item \textbf{RLHF（Reinforcement Learning from Human Feedback）}：reward 来自人类偏好（或偏好模型），目标是让模型输出更符合人类期望。适用于聊天、创意写作等主观任务。
\item \textbf{RLVR（Reinforcement Learning from Verifiable Rewards）}：reward 是可客观验证的（数学题答案、代码是否通过测试、Agent 任务是否达成目标）。DeepSeek V3 的成功很大程度上证明了 RLVR 的价值——直接从 pretrain checkpoint 出发，通过 RLVR 就能大幅提升模型智能。
\end{itemize}

\begin{importantbox}{RLVR 的范式意义}
DeepSeek V3 的 release 是一个转折点。在此之前，很多人以为「做好 SFT 就够了」——直接从 GPT 等模型的输出做蒸馏就能得到好的开源模型。DeepSeek V3 证明了 RLVR 能直接提升模型的 intelligence，甚至从 pretrain checkpoint 出发就足够。这件事引起了包括 NVIDIA 在内的各方高度关注，成为推动 RL 在业界普及的关键事件。
\end{importantbox}

\subsection{本章小结}
强化学习在 LLM 训练中的地位已经从「锦上添花」升级为「核心能力提升引擎」。PPO 上限高但工程难度大（核心是 Critic 的初始化），GRPO 更简单但可能牺牲上限，RLVR 在可验证领域展现了最强的 scaling potential。数据仍然是最大的瓶颈——好的偏好数据和 RL 环境数据是稀缺资源。

\section{NexusFlow、英伟达与二次创业}

\subsection{NexusFlow：早期 Agent 的先驱}

NexusFlow 的创业方向是 Agent——让模型能主动使用工具（tool use），在多轮交互中自主完成复杂任务。朱邦华回忆，在 2023--2024 年：

\begin{quote}
「大家会问什么是 Agent？是不是就是 RAG——retrieval augmented generation？是不是就是搜索一下文本然后给个 response？我们需要花很多时间去教育市场。」\par
\hfill —— 朱邦华，00:41:00--00:41:30
\end{quote}

虽然团队比较早地预判了 Agent 的方向，但市场教育的成本比预期更高，Agent 真正爆发的时间点也比预期晚了一些。

\subsection{英伟达的体验：NEOJIM 项目}

被英伟达收购后，朱邦华发起了一个名为 NEOJIM 的内部项目——相当于英伟达内部的「GPT 计划」：

\begin{quote}
「NVIDIA 是一个非常大的组织，每个地方都有各种好的想法和好的 data。我们希望把所有内部的 Agentic data 收集起来，用在自己模型的训练里。」\par
\hfill —— 朱邦华，00:42:30--00:43:00
\end{quote}

这个项目从一个小 team 的自发努力开始，后来逐渐升级为公司级战略。朱邦华在英伟达体验到了「大厂的资源密度」——当他需要千卡甚至万卡级别的 GPU 来做实验时，这在创业公司是不可想象的。

\begin{knowledgebox}{千卡级别的大规模训练挑战}
朱邦华指出，单机或几个节点的训练相对简单，但一旦到了千卡甚至万卡级别，系统 Infra 的要求完全不同：
\begin{itemize}
\item 需要做 fault tolerance（故障容错）
\item 需要在训练速度和稳定性之间做 tradeoff
\item 调试难度指数级上升——「你浪费很多各种各样不同 case 的 run，才能理解这个东西到底是否真的 reliable」
\end{itemize}
这也是为什么他认为「卡多」本身就是一个核心竞争壁垒——它能让你在更高维度上试错和学习。
\end{knowledgebox}

\subsection{不到半年就离开：勇于放弃}

最令人震撼的是，朱邦华在英伟达被收购后不到半年就选择离开——这意味着放弃后续大量的归属权益。当被问及放弃了多少钱时，他笑着说「具体不太能说，但可能也算是足够躺平生活的」，随即补了一句「肯定不如我们影老板（盛颖）放弃的多」。

\begin{quote}
「我老婆和父母都非常支持。」\par
\hfill —— 朱邦华，00:47:30--00:47:45
\end{quote}

他和盛颖决定以 SGLang 为基础，重新出发。RadixArk 很快拿到了 1 亿美元融资，估值超过 4 亿美元。

\subsection{本章小结}
NexusFlow 的早期探索让朱邦华积累了 Agent 方向的实战经验和对市场的判断力。英伟达的经历则让他见识了大厂级别的资源密度和大规模 Infra 挑战。从被收购到再次创业，转折的核心驱动力是：他看到了以 SGLang 为基础的更大机会，而且他愿意为了这个机会放弃已经到手的巨额财富。


\section{AI Infra 的大规模挑战}

\subsection{从单卡到万卡：问题的质变}

朱邦华用一个生动的类比区分了不同规模下的挑战：

\begin{itemize}
\item \textbf{单机/几个节点}：系统搭起来比较容易，即使训练一个 100B 级别的模型也不会特别复杂
\item \textbf{千卡/万卡级别}：问题的本质发生变化——fault tolerance 不再是可选项，而是必选项；调度策略、通信拓扑、梯度同步的频率和精度都需要重新设计
\end{itemize}

\begin{importantbox}{大规模 RL 训练的独特挑战}
强化学习训练与预训练在 Infra 需求上有根本不同：预训练是相对确定性的流水线（forward + backward），而 RL 训练涉及多个模型（Actor、Critic、Reward Model）的交替更新，以及 rollout 的生成与评估。在万卡规模下协调这些步骤，需要全新的调度和编排系统。
\end{importantbox}

\subsection{直觉驱动的实验管理}

朱邦华强调了大规模训练中的一个「反直觉」洞见：不要过度优化。

\begin{quote}
「我觉得这个 run 它就是好的，我就直接 launch。我不会说我一定要精打细算、做很多 ablation，因为我靠 intuition 来做一些比较正确的选择，在非常 high dimension 的 search space 里面去寻找最好的 data、算法和 Infra 的 combination。」\par
\hfill —— 朱邦华，00:46:30--00:47:05
\end{quote}

这个观点背后是一个更深的工程哲学：在大规模训练中，搜索空间极其高维，穷举式的 ablation study 是资源浪费。真正高效的做法是：依靠经验和直觉快速迭代，在「差不多对」的时候果断 launch，从实际结果中学习。

\subsection{本章小结}
大规模训练的挑战不是线性的——它会在千卡级别发生质变。核心矛盾在于：高维搜索空间 vs 有限计算资源。朱邦华的解决方式是「直觉 + 快速迭代」，而非过度优化的 ablation study。这种思路与他在学术和创业中展现的判断力一脉相承。

\section{中美 AI 生态对比}

\subsection{开源大模型：差距正在缩小}

朱邦华对中美开源大模型的评估相当坦率：

\begin{itemize}
\item \textbf{Infra 能力}：国内团队的 Infra 水平很高，很 solid
\item \textbf{数据}：这是国内主要的受限方向——高质量数据（尤其是复杂的 RL 环境数据）比较稀缺
\item \textbf{算法}：中美差距不大，甚至在快速追赶
\end{itemize}

\begin{quote}
「我觉得这个差距是在慢慢缩短。美国的开源大模型和中国的开源大模型差距会越来越小。」\par
\hfill —— 朱邦华，00:24:00--00:24:45
\end{quote}

\subsection{中国工程师的独特优势}

当主持人提到硅谷流传的一个说法——「如果 AI Infra 项目里中国人数量不能超过一半，代码质量会急剧下降」——朱邦华给出了一个微妙的回应：

\begin{quote}
「美国人也有很多非常强的。但我觉得现在这个时代，AI 主要看 coding 和一些数学能力。国内同学从 1 做到 100，做到 perfect 的能力是非常非常强的。」\par
\hfill —— 朱邦华，00:23:00--00:23:45
\end{quote}

\begin{knowledgebox}{「从 1 到 100」vs「从 0 到 1」}
朱邦华区分了两种能力：（1）探索阶段（从 0 到 1）——需要大量试错，很多尝试不会 work，但偶尔有突破就能创造新东西；（2）工程阶段（从 1 到 100）——需要把已知可行的事情做到极致。他认为国内工程师在后者上的能力是毋庸置疑的，这也是 SGLang 社区能迅速壮大的原因之一。
\end{knowledgebox}

\subsection{本章小结}
中美 AI Infra 的差距在快速缩小。中国团队的工程落地能力极强（「从 1 到 100」），但原始创新（「从 0 到 1」）和数据积累仍是短板。SGLang 作为一个中美混合团队（伯克利起源 + 中国工程力量）的成功，恰恰证明了两种优势互补的可能性。

\section{AI 重塑计算机行业：Agent 时代的来临}

\subsection{程序员会被取代吗？}

这是访谈中最具争议也最引人深思的话题。朱邦华的观点非常直白：

\begin{quote}
「Agent 可以取代一些更 basic 的 CS/software 相关的工作。另一方面，它会让很多资源更 concentrate 到真正强和有 taste 的、有水平的顶尖 system builder 或者 software engineer 上，甚至是 researcher 上。」\par
\hfill —— 朱邦华，00:93:30--00:94:15
\end{quote}

他的核心逻辑是：Agent 会无限放大强者的能力，加速他们的 iteration speed，从而导致：
\begin{itemize}
\item 以前需要一百人维护的复杂软件系统，可能只需要一个小团队 + Agent
\item 中间层工程师的需求会急剧萎缩
\item 顶尖工程师和研究者的价值反而会被放大
\end{itemize}

\begin{warningbox}{CS 毕业生的困境}
朱邦华坦承他也不知道答案：「计算机是一个非常大的专业，基本上每个学校计算机的人数都能排到前几名。接下来的话，真的看不太清楚这么多计算机专业毕业生怎么继续在这个行业里生存。」
\end{warningbox}

\subsection{不只是危机，也是解放}

但朱邦华也给出了一个更乐观的视角：

\begin{quote}
「如果有些人就不喜欢 CS，但因为 AI 的原因，他其实是可以被 empower 去做任何他想做的事情。这个换一个角度来说，其实也是好事。」\par
\hfill —— 朱邦华，00:97:00--00:97:30
\end{quote}

他区分了两类人：真正享受 coding 的人可以「古法编程」追求兴趣和审美，不享受的人则可以被 AI 解放去做其他事情。这个观点与马斯克「物质极大自由后，大家做自己开心的事」的愿景有异曲同工之处——但朱邦华也清醒地指出：「任何社会的转型都会带来很多阵痛。有些人就是享受写 code，但能力不是那么强，这种人确实会面临找不到工作的挑战。」

\subsection{本章小结}
AI Agent 正在从根本上重塑软件工程的劳动力结构：中间层被压缩，顶层被放大。这不是一个「会不会发生」的问题，而是「如何适应」的问题。朱邦华的建议不是盲目乐观，而是鼓励找到自己真正热爱的东西——因为热爱是 AI 无法替代的驱动力。

\section{判断力、品味与高能动性个体}

\subsection{判断力从何而来？}

当主持人问朱邦华「你的判断力为什么这么准」时，他没有给出神秘化的解释，而是将它归因于大量的探索和试错：

\begin{quote}
「我认识的很多人，大家其实都是不断去探索，然后在自己慢慢探索到一些比较好的机会的时候，那一次或那一篇工作让自己拔高了高度。但你也需要判断什么时候是版本到来的时候。」\par
\hfill —— 朱邦华，00:64:30--00:65:15
\end{quote}

他用游戏术语「版本砸脸」来描述机遇——时代会提供窗口，但能否识别并抓住窗口，取决于之前的积累和当下的判断。

\subsection{品味：比智商更稀缺的能力}

在访谈的结尾部分，朱邦华和主持人讨论了一个核心概念——「品味」（taste）：

\begin{quote}
「说实话，你要说智商超过 180 或 200 的人，几乎也见不到。但大多数在这种细小的差别之下，全靠品味——这个事怎么就不美呢？这个事做得怎么就不漂亮呢？这个事我为什么能用 500 行代码写出来，能不能用 100 行写出来，而且大家看得更懂？」\par
\hfill —— 朱邦华，00:102:30--00:103:15
\end{quote}

\begin{importantbox}{品味的三层含义}
朱邦华和主持人讨论的「品味」可分解为三层：
\begin{enumerate}
\item \textbf{审美层}：看一段代码，能本能地感觉到「不对劲」——这不是功能上的 bug，而是结构上的丑陋
\item \textbf{简化层}：能用 100 行写出别人需要 500 行的功能，且更清晰易懂。这不是炫技，而是对本质的把握
\item \textbf{判断层}：在众多可能的方向中，直觉上知道哪个是对的——不是通过 elimination，而是通过更深的 pattern recognition
\end{enumerate}
\end{importantbox}

\subsection{高能动性个体}

朱邦华提到了「high-agency individual」这个概念。他以 Linus Torvalds 为例——Linus 在 Linux 内核开发中的做法，本质上就是早期「用 Agent 完成工作」的雏形：把代码审查的标准提升到极致，任何不符合审美的代码直接 reject，不管贡献者是谁。

\begin{quote}
「就是这个东西到底有什么问题，这个其实就是你在面试或者交流的时候能感受到的——真的有人聊的时候就说，我现在看公司里的代码浑身不自在，我就希望把所有不好的代码都删掉。」
\end{quote}

\subsection{本章小结}
判断力来自大量探索后的模式识别，品味是一种对「什么是对的」的本能感知，高能动性则是将判断和品味转化为行动的能力。这三者在 AI 时代形成了正反馈：AI 放大强者的能力，而判断力、品味和能动性恰恰是 AI 最难以复制的品质。

\section{未来愿景与给年轻人的建议}

\subsection{个人 AI 代理：RadixArk 的长远愿景}

当谈到 RadixArk 的终极愿景时，朱邦华给出了一个宏大的画面：

\begin{quote}
「会不会有一家公司需要去 empower 所有的 personal agent？这些 agent 之间的交互、以及和人交互的 interface，是不是都会需要一家公司来定义？」\par
\hfill —— 朱邦华，00:77:00--00:77:45
\end{quote}

他的比喻是：就像现在人用电一样自然——RadixArk 想做的是那个「AI 电厂」，让每个人都能拥有自己的 AI Agent。他同时坦承这「只是万千未来可能性中的一种」，也有可能 OpenAI 或 Anthropic 一家独大，把游戏结束。但他更希望看到百花齐放的时代。

\subsection{Neo-Labs：鼓励探索的生态}

朱邦华对 Neo-Labs（新时代的 AI 研究 Lab）持强烈支持态度：

\begin{quote}
「你就算这次不成，你可能还会有下次。但如果你一直不去做的话，你可能就没有机会会成。」\par
\hfill —— 朱邦华，00:83:00--00:83:30
\end{quote}

他认为 OpenAI 和 Anthropic 逼近万亿市值的示范效应已经足够明显：如果你投入几亿或十几亿，万一出现一次范式革命，很可能「1 亿就变成 1 万亿」。所以无论成与不成，都应该去做——探索本身就是价值。

\subsection{给三类人的建议}

在访谈最后，朱邦华给出了他给三类人的真诚建议：

\textbf{给本科生}：广泛探索，不要过早给自己设限。「我自己在本科的时候也不知道以后会做什么。先去探索，找到自己真正感兴趣的东西。」

\textbf{给博士生}：如果发现自己真正想做的事情不在当前轨道上，勇于调整。「为什么这个人他就能做这个，为什么我就做不了？你试一下，说不定你做的其实比他还好。」

\textbf{给工作的人}：保持对「美」的追求。如果看公司里的代码浑身不自在，那说明你的品味在起作用——不要因为环境而降低标准，去找到那些和你一样「有强迫症」的人。

\begin{importantbox}{一个贯穿始终的建议}
朱邦华最核心的建议可以浓缩为一句话：「如果你看到有人做了一件事，心里想『他凭什么』或者『为什么我不行』——那就真的去试一下。很可能你做了之后发现，这个人其实也没什么大不了的。」
\end{importantbox}

\subsection{本章小结}
朱邦华的未来愿景是「人人都有 AI Agent」的百花齐放时代，他鼓励探索（Neo-Labs）、支持试错（不成还有下次），并给三类年轻人提供了具体而真诚的建议。而他最核心的哲学——勇于放弃、敢于尝试——贯穿了整个人生轨迹。


\section{总结与延伸}

\subsection{讲者的核心洞见}

回顾这场 107 分钟的对话，朱邦华的核心洞见可以归纳为以下几个层次：

\textbf{技术层}
\begin{itemize}
\item Post-training（尤其是 RL）是当前 AI 基础设施最大的机遇——比推理更被低估，比预训练更具普惠性
\item 推理引擎的核心竞争力从「通用性能」转向「架构适配速度」——谁能最快支持新模型架构，谁就赢
\item PPO 训练的核心工程难点在于 Critic 的初始化——从 Reward Model 热启动可以大幅提升成功率
\item RLVR 在可验证领域展现了最强的 scaling potential，但复杂的 RL 环境数据仍是瓶颈
\end{itemize}

\textbf{创业层}
\begin{itemize}
\item 「勇于放弃」不是鲁莽，而是基于判断力的理性选择——看到更大的机会时，果断转向
\item 家庭支持是创业者的底层操作系统——朱邦华三次提到家人的支持
\item 不要等「完美的 run」，靠直觉在正确的时候 launch——高维搜索空间中，ablation study 不如快速迭代
\end{itemize}

\textbf{人生哲学层}
\begin{itemize}
\item 判断力来自大量探索 + 模式识别——没有捷径，但有迹可循
\item 品味比聪明更重要——能感受「什么是美」「什么不对劲」是 AI 时代最稀缺的能力
\item 「他凭什么」的态度是创业者的第一驱动力——不是嫉妒，是自信
\end{itemize}

\subsection{综合与延伸}

\begin{figure}[H]
\centering
\includegraphics[width=0.75\textwidth]{figures/talk_closing.jpg}
\caption{访谈临近结束：朱邦华分享给年轻人的建议 \protect\footnotemark}
\end{figure}
\footnotetext{视频画面时间区间：01:44:00--01:45:15。}

朱邦华的职业路径展示了一种罕见的「判断力密度」：

\begin{enumerate}
\item \textbf{第一次判断}（2023 年初）：GPT 是范式转变，不是炒作——从理论转向系统和 LLM
\item \textbf{第二次判断}（2023 年中）：PPO + RLHF 是大模型能力提升的必经之路——创办 NexusFlow 率先突破
\item \textbf{第三次判断}（2024 年底）：英伟达不是终点，SGLang 的社区和生态有更大的想象空间——放弃股权，二次创业
\end{enumerate}

每一次判断都涉及「放弃」——放弃已经积累的学术资本、放弃大厂的稳定和高薪、放弃到手的巨额财富。但这种放弃不是盲目的，而是基于一个更深的信念：AI 浪潮才刚刚开始，早期所有人都有机会。

\begin{importantbox}{「111」公式：朱邦华的反事实思维}
朱邦华反复使用的一种思维方式可以概括为「111」：
\begin{itemize}
\item 1 亿投入 → 1 万亿可能 → 1 次范式革命
\end{itemize}
这个公式不是财务分析，而是一种风险认知：在 AI 的早期阶段，不做的成本（错失范式革命）远大于做的成本（亏掉几亿）。这种思维正是他所有「勇于放弃」决策的底层逻辑。
\end{importantbox}

\subsection{交叉连接}

访谈中闪现的几条线索之间存在深刻的关联：

\begin{itemize}
\item \textbf{「Long context」能力与判断力}：朱邦华在清华展现的快速吸收信息的能力，与他后来能准确判断 AI 趋势的能力本质上是同一种认知模式——在有限时间内，从大量信息中提取关键信号
\item \textbf{Michael Jordan 的经济学训练与 RadixArk 的商业模式}：Jordan 关于「Human-AI interaction 本质上是经济学问题」的洞见，隐隐指向了 RadixArk 想做「Agent-to-Agent 交互平台」的终极愿景
\item \textbf{PPO 的 Critic trick 与创业哲学}：Critic 热启动的本质是「不要从零开始，要站在已有的认知上迭代」——这与他不盲目推倒重来、而是在 SGLang 已有社区基础上创业的思路完全一致
\end{itemize}

\begin{warningbox}{Chatbot Arena 的局限性}
朱邦华对 LMSYS Chatbot Arena 的评价值得注意：Arena 的评分主要反映「用户喜不喜欢这个回复」，而当前的前沿需求是「模型能不能解决最复杂的 coding 问题」——这两者之间没有必然联系。过度依赖 Arena 排名可能导致对模型实际 intelligence 的误判。这个观点对正在追榜的 AI 团队是一个重要的提醒。
\end{warningbox}

\subsection{开放问题}

朱邦华在访谈中含蓄地提出了几个尚未解决的问题：

\begin{enumerate}
\item \textbf{智能爆炸的拐点}：模型能力会不会出现 exponential 的增长拐点？如果出现，谁会先到达？（他的态度是「希望相信，但受制于能源等现实约束」）
\item \textbf{Interaction Model 的未来}：下一代交互模型（如 DeepSeek V4 的新交互范式）会如何改变推理 Infra 的设计？（他暗示有重大合作在进行中，但暂时无法公开）
\item \textbf{CS 专业的教育改革}：当 AI Agent 能完成大部分 basic coding 时，大学计算机专业应该教什么？（他坦承「不知道，这是一个很难的问题」）
\item \textbf{Agent 之间的交互协议}：如果未来每个人都有自己的 AI Agent，Agent 之间如何通信、协调、竞争？（这是他给 RadixArk 设定的长远研究课题）
\end{enumerate}

\subsection{拓展阅读}

\begin{itemize}
\item SGLang 项目主页：\href{https://github.com/sgl-project/sglang}{github.com/sgl-project/sglang} —— 开源 LLM 推理引擎
\item FlexGen 论文：\href{https://arxiv.org/abs/2303.06865}{FlexGen: High-Throughput Generative Inference of Large Language Models with a Single GPU} —— SGLang 的前身之一
\item DeepSeek V3 技术报告 —— 展示了 RLVR 在大规模模型训练中的有效性
\item RadixArk 官方网站 —— 公司最新动态和产品信息
\end{itemize}

%% --- 正文内容结束 --- %%

\end{document}
